% Slideaway — Conference Talk Template (15–20 min)
% Theme: Madrid/beaver — classic institutional style for AGU/EGU/AMS/NeurIPS
% Structural scaffolding only. Style/color theory → design-foundations.md
% Engine workflow → engine-beamer.md

\documentclass[aspectratio=169,11pt]{beamer}

\usetheme{Madrid}
\usecolortheme{beaver}

% Remove navigation clutter — frame numbers suffice for conference talks
\setbeamertemplate{navigation symbols}{}
\setbeamertemplate{footline}[frame number]

% ── Packages ──────────────────────────────────────────────────────────────
\usepackage{booktabs}
\usepackage{graphicx}
\usepackage{amsmath}
\usepackage{hyperref}
\usepackage{xcolor}
\usepackage[backend=biber,style=authoryear]{biblatex}
% \addbibresource{references.bib}  % Uncomment and point to your .bib file

% ── Title Metadata ────────────────────────────────────────────────────────
\title[Short Title]{Full Presentation Title}
\subtitle{Conference Name — Session ID}
\author[A.~Author]{Author Name\inst{1} \and Co-Author Name\inst{2}}
\institute[Inst]{
  \inst{1}Department of X, University of Y \and
  \inst{2}Agency Z
}
\date{\today}

% ══════════════════════════════════════════════════════════════════════════
\begin{document}

\begin{frame}
  \titlepage
\end{frame}

% ── Introduction (2–3 slides, ~3 min) ─────────────────────────────────────
\section{Introduction}

\begin{frame}{The Problem}
  % Hook: surprising statistic, provocative question, or data teaser
  \begin{itemize}
    \item<1-> Context and motivation — why this matters now
    \item<2-> Background — what is known
    \item<3-> Research question or hypothesis
  \end{itemize}
  \note{Open with energy. State the gap clearly. Keep background to 90 seconds max.}
\end{frame}

% ── Methods (2 slides, ~3 min) ────────────────────────────────────────────
\section{Methods}

\begin{frame}{Approach}
  \begin{columns}[T]
    \begin{column}{0.48\textwidth}
      \textbf{Design}
      \begin{itemize}
        \item Data sources and coverage
        \item Experimental design or model setup
        \item Key parameters and choices
      \end{itemize}
    \end{column}
    \begin{column}{0.48\textwidth}
      % Replace with your method diagram
      \framebox[\textwidth]{\parbox{0.9\textwidth}{\centering\vspace{2cm}Method Diagram\vspace{2cm}}}
    \end{column}
  \end{columns}
  \note{Emphasize what is novel. Skip standard details the audience can read in the paper.}
\end{frame}

\begin{frame}{Analysis Pipeline}
  \begin{enumerate}
    \item Pre-processing and quality control
    \item Core analysis or model execution
    \item Post-processing and evaluation metrics
  \end{enumerate}
  \note{Keep it high-level. Backup slides cover implementation details.}
\end{frame}

% ── Results (6–8 slides, ~8 min) ──────────────────────────────────────────
\section{Results}

\begin{frame}{Result 1 — Main Finding}
  \begin{center}
    % Replace with your key figure
    \framebox[0.8\textwidth]{\parbox{0.7\textwidth}{\centering\vspace{3cm}Key Result Figure\vspace{3cm}}}
  \end{center}
  \note{Lead with the figure. State the takeaway in one sentence. Let the visual carry the weight.}
\end{frame}

\begin{frame}{Result 2 — Supporting Evidence}
  \begin{columns}[T]
    \begin{column}{0.48\textwidth}
      \framebox[\textwidth]{\parbox{0.9\textwidth}{\centering\vspace{2cm}Figure A\vspace{2cm}}}
    \end{column}
    \begin{column}{0.48\textwidth}
      \framebox[\textwidth]{\parbox{0.9\textwidth}{\centering\vspace{2cm}Figure B\vspace{2cm}}}
    \end{column}
  \end{columns}
  \note{Compare side-by-side. Point out the specific feature that supports your argument.}
\end{frame}

\begin{frame}{Result 3 — Quantitative Summary}
  \begin{table}
    \caption{Performance comparison}
    \begin{tabular}{lcc}
      \toprule
      Method & Metric A & Metric B \\
      \midrule
      Baseline & 0.72 & 0.81 \\
      Proposed & \textbf{0.89} & \textbf{0.94} \\
      \bottomrule
    \end{tabular}
  \end{table}
  \note{Bold your best numbers. Mention statistical significance if applicable.}
\end{frame}

% ── Discussion (1–2 slides, ~2 min) ───────────────────────────────────────
\section{Discussion}

\begin{frame}{Interpretation and Implications}
  \begin{itemize}
    \item What the results mean in context of existing work
    \item Broader implications for the field
    \item Known limitations and caveats
  \end{itemize}
  \note{Be honest about limitations. Reviewers and audiences respect candor.}
\end{frame}

% ── Conclusion (1 slide, ~1 min) ──────────────────────────────────────────
\section{Conclusion}

\begin{frame}{Key Takeaways}
  \begin{enumerate}
    \item First main conclusion
    \item Second main conclusion
    \item Future work direction
  \end{enumerate}
  \vfill
  \small
  Code: \url{https://github.com/your/repo} \\
  Contact: \href{mailto:you@example.com}{you@example.com}
  \note{End with 3 crisp assertions. Show the repo link. Invite questions.}
\end{frame}

% ── Acknowledgments ───────────────────────────────────────────────────────

\begin{frame}{Acknowledgments}
  \small
  This work was supported by [Funding Agency, Grant \#].

  We thank [collaborators/reviewers] for helpful discussions.

  % \printbibliography  % Uncomment if using biblatex
\end{frame}

% ── Q\&A ──────────────────────────────────────────────────────────────────

\begin{frame}[standout]
  Questions?
\end{frame}

% ══════════════════════════════════════════════════════════════════════════
\appendix

\begin{frame}{Backup: Detailed Methods}
  % Extended methodology for anticipated questions
  \note{Have 3–5 backup slides ready. Label them clearly.}
\end{frame}

\begin{frame}{Backup: Sensitivity Analysis}
  % Parameter sensitivity, robustness checks
\end{frame}

\begin{frame}{Backup: Additional Results}
  % Supplementary figures/tables that didn't fit the main talk
\end{frame}

\end{document}
